\chapter{Valutazione e Ottimizzazione Ambientale del Livello Front-End}

L'analisi del Front-End del sistema 2Chance richiede una calibrazione concettuale significativa rispetto ad altre architetture web moderne. Mentre i framework Single Page Application (SPA) contemporanei, come React utilizzato nel progetto di riferimento Rently, scaricano la quasi totalità del carico computazionale di rendering del Document Object Model (DOM) sul browser client (sfruttando ed esaurendo direttamente la batteria dello smartphone o del laptop dell'utente finale), l'architettura di 2Chance è ancorata al Server-Side Rendering tramite JavaServer Pages (JSP).

Nel modello adottato da 2Chance, è il container servlet Tomcat a processare la logica visiva, interpolare i dati estratti dai Model Beans e assemblare iterativamente le stringhe HTML prima di trasmetterle attraverso la rete. Questa differenza architetturale non annulla la questione ecologica, ma la trasla nello spazio. L'elaborazione lato server sposta il consumo energetico verso i data center (che teoricamente potrebbero essere alimentati da fonti rinnovabili), ma produce payload di testo HTML completi che risultano significativamente più pesanti, in termini di byte, rispetto alle agili risposte in formato JSON consumate dalle API REST delle architetture React. Il trasferimento di questi byte aggiuntivi attraverso l'infrastruttura di rete globale (router, switch, ponti radio, reti cellulari) genera un impatto ecologico massiccio e spesso sottostimato.

\section{Analisi Ambientale delle Pagine JSP (GreenIT ed EcoIndex)}
